\begin{proposition}[Killing 方程的协变表达]
在 Riemann 流形上，\[\mathcal L_Xg=0\quad\Longleftrightarrow\quad\nabla_iX_j+\nabla_jX_i=0.\]对 Killing 场，仿射参数测地线上的 $g(X,\dot\gamma)$ 恒定。
\end{proposition}

\begin{proof}
将 $\mathcal L_Xg$ 的定义用 Levi--Civita 联络改写，利用度量相容性和无挠性消去高阶项，得到一阶方程。

对任意向量场 $Y,Z$，Lie 导数定义给出
\begin{equation*}
 (\mathcal L_Xg)(Y,Z)
 =X[g(Y,Z)]-g([X,Y],Z)-g(Y,[X,Z]).
\end{equation*}
这是 Lie 导数的 Cartan 公式在 $(0,2)$ 张量上的特例：第一项是 $X$ 作用于张量值，后两项扣除 $Y,Z$ 自身被 $X$ 的流推动产生的变化。

Levi--Civita 联络满足 $\nabla g=0$，故
\begin{equation*}
 X[g(Y,Z)]
 =(\nabla_Xg)(Y,Z)+g(\nabla_XY,Z)+g(Y,\nabla_XZ)
 =g(\nabla_XY,Z)+g(Y,\nabla_XZ).
\end{equation*}
度量相容性 $\nabla g=0$ 是 Levi--Civita 联络的两条公理之一，几何意义是"协变微分保持内积"。

$\nabla$ 无挠给出 $[X,Y]=\nabla_XY-\nabla_YX$、$[X,Z]=\nabla_XZ-\nabla_ZX$。代入上面的 Lie 导数定义，
\begin{align*}
 (\mathcal L_Xg)(Y,Z)
 &=g(\nabla_XY,Z)+g(Y,\nabla_XZ)\\
 &\quad -g(\nabla_XY-\nabla_YX,Z)-g(Y,\nabla_XZ-\nabla_ZX).
\end{align*}
无挠性是 Levi--Civita 联络的另一条公理，把 Lie 括号这一纯微分拓扑对象用协变导数表达。

含 $\nabla_XY$ 与 $\nabla_XZ$ 的项相消，留下
\begin{equation*}
 (\mathcal L_Xg)(Y,Z)
 =g(\nabla_YX,Z)+g(Y,\nabla_ZX).
\end{equation*}
此式对任意 $Y,Z$ 成立。注意 $X$ 现在被 $\nabla$ 作用，而不是作用于别人——这是 Killing 方程为一阶 PDE 的关键。

取 $Y=\partial_i$、$Z=\partial_j$，$g(\partial_i,\partial_j)=g_{ij}$，$\nabla_{\partial_i}X$ 的 $j$ 分量为 $\nabla_iX^j$，故
\begin{equation*}
 (\mathcal L_Xg)_{ij}
 =g_{jk}\nabla_iX^k+g_{ik}\nabla_jX^k
 =\nabla_iX_j+\nabla_jX_i.
\end{equation*}
$\mathcal L_Xg=0$ 即给出
\begin{equation}\label{eq:math-dg-synthesis-killing-eq}
{\nabla_iX_j+\nabla_jX_i=0,}
\end{equation}
即\eqnrefs{eq:math-dg-synthesis-killing-eq}。这是一阶线性 PDE。
沿测地线计算 $\nabla_TT=0$ 及 Killing 方程，得到
\[
\frac{\dd}{\dd s}g(X,T)=g(\nabla_TX,T)=0.
\]
Killing 场的限制是测地线变分的 Jacobi 场。
沿每条径向测地线，$X$ 由 $X(p)$ 与 $\nabla X(p)$ 的初值唯一决定；
Killing 方程又要求 $\nabla X(p)$ 反对称。
故在连通流形上，其第一 jet 的可能维数至多为
$n+\binom n2=n(n+1)/2$，零初值由 Jacobi ODE 的唯一性给出零场。



典型例子。
\begin{itemize}
\item $\RR^n$：平移 $\partial_i$ 与旋转 $x^i\partial_j-x^j\partial_i$ 共 $n+n(n-1)/2=n(n+1)/2$ 个 Killing 场，达到最大维数。
\item $S^n$：$SO(n+1)$ 的 Lie 代数给出 $n(n+1)/2$ 个 Killing 场，对应球面的等距群。
\item $\HH^n$：$SO(n,1)$ 给出 $n(n+1)/2$ 个 Killing 场，双曲空间也是最大对称空间。
\item 一般 Riemann 流形：Killing 场维数 $<n(n+1)/2$，反映对称性的缺失。
\end{itemize}

Killing 场对应度量的连续对称，沿测地线产生守恒量 $J_X=g(X,\dot\gamma)$。在广义相对论中，Killing 场对应时空对称（如稳态时空的时间平移 $\partial_t$ 给出能量守恒、轴对称 $\partial_\phi$ 给出角动量守恒）。Noether 定理在几何语言下即"Killing 场 $\Rightarrow$ 测地线守恒量"。在经典力学中，$n(n+1)/2$ 维最大对称空间对应 $n$ 维自由粒子全部运动积分的存在性。
\end{proof}

\begin{proposition}[Hodge 拉普拉斯的椭圆主符号]
设 $\Delta_H=\dd\delta+\delta\dd$ 作用于 Riemann 流形的 $k$-形式。采用 $\partial_j\mapsto\ii\xi_j$，\[\sigma_2(\Delta_H)(\xi)=|\xi|_g^2\operatorname{id}_{\Lambda^kT_p^*M}.\]故 $\Delta_H$ 椭圆。
\end{proposition}

\begin{proof}
计算 $\Delta=\dd\delta+\delta\dd$ 的二阶主符号。$\dd$ 与 $\delta$ 的一阶主符号分别由外积与插入算子给出，二者复合后用 Cartan 恒等式 $\iota_{\xi^\sharp}(\xi\wedge\cdot)+\xi\wedge\iota_{\xi^\sharp}\cdot=|\xi|^2\,\cdot$ 化为 $|\xi|^2\id()$。

在点 $p$ 处取非零余向量 $\xi\in T_p^*M$。采用 $\partial_j\mapsto \ii\xi_j$ 的主符号约定，外微分与余微分的一阶主符号分别为
\begin{equation*}
\sigma_1(\dd)(\xi)=\ii\,\xi\wedge,
\qquad
\sigma_1(\delta)(\xi)=-\ii\,\iota_{\xi^\sharp}.
\end{equation*}
$\sigma_1(\dd)=\ii\,\xi\wedge$ 来自外微分对一形式 $\alpha$ 给 $\dd\alpha=\partial_i\alpha_j\,\dd x^i\wedge\dd x^j$，主项含一阶导数。$\sigma_1(\delta)=-\ii\,\iota_{\xi^\sharp}$ 来自余微分的 Hodge 星定义与 $\delta=\pm *\dd*$ 中 $*$ 不改变阶数。

因此二阶算子 $\Delta=\dd\delta+\delta\dd$ 的主符号为
\begin{equation*}
\begin{aligned}
\sigma_2(\Delta)(\xi)
&=\sigma_1(\dd)\sigma_1(\delta)+\sigma_1(\delta)\sigma_1(\dd)\\
&=(\ii\,\xi\wedge)(-\ii\,\iota_{\xi^\sharp})
+(-\ii\,\iota_{\xi^\sharp})(\ii\,\xi\wedge)\\
&=(\xi\wedge)\iota_{\xi^\sharp}
+\iota_{\xi^\sharp}(\xi\wedge).
\end{aligned}
\end{equation*}
两个 $\ii$ 与 $-\ii$ 相乘都给 $+1$，符号约定使最终结果为正。

对任意形式 $\alpha$，插入算子与外积满足
\begin{equation*}
\iota_{\xi^\sharp}(\xi\wedge\alpha)
=|\xi|_g^2\alpha
-\xi\wedge\iota_{\xi^\sharp}\alpha.
\end{equation*}
这是 Cartan 公式 $\mathcal L_{\xi^\sharp}=\dd\iota_{\xi^\sharp}+\iota_{\xi^\sharp}\dd$ 在常向量场上的特例，也可直接对单项式 $\alpha=\theta^{i_1}\wedge\cdots\wedge\theta^{i_k}$ 验证：若 $\xi$ 出现在 $\alpha$ 中，插入与外积相继给出 $|\xi|^2$；若不出现，外积先添加该方向，插入再移除，仍给出 $|\xi|^2\alpha$；两种情形都是同一个反对易恒等式。

把 Cartan 恒等式代入主符号表达式，
\begin{equation*}
\iota_{\xi^\sharp}(\xi\wedge\alpha)
+\xi\wedge\iota_{\xi^\sharp}\alpha
=|\xi|_g^2\alpha
-\xi\wedge\iota_{\xi^\sharp}\alpha
+\xi\wedge\iota_{\xi^\sharp}\alpha
=|\xi|_g^2\alpha,
\end{equation*}
两项相消，恰好得到
\begin{equation}\label{eq:math-dg-hodge-principal-symbol}
{
\sigma_2(\Delta)(\xi)
=|\xi|_g^2\id(\Lambda^kT_p^*M).}
\end{equation}
当 $\xi\neq0$ 时该线性映射可逆，故 $\Delta$ 是椭圆算子。结合\eqnrefs{eq:math-dg-hodge-energy}的非负性与适当自伴定义域，就能调用前面建立的自伴谱理论，而不必把 Hodge 分解当作孤立结论。

典型例子与对比。
\begin{itemize}
\item $\RR^n$ 上 $0$-形式（函数）：$\Delta f=-\sum_i\partial_i^2 f$（本书非负约定），主符号 $|\xi|^2$，正是普通 拉普拉斯算子。
\item $\RR^n$ 上 $1$-形式：$\Delta\alpha_j=-\sum_i\partial_i^2\alpha_j$，每个分量独立满足标量 拉普拉斯算子，主符号仍是 $|\xi|^2\id()$。
\item 一般流形：主符号与 $\RR^n$ 相同，但低阶项含曲率（见 Weitzenböck 公式 $\Delta_H=\nabla^*\nabla+\operatorname{Ric}$），曲率通过低阶项影响谱而非椭圆性。
\end{itemize}

Hodge 拉普拉斯算子 的椭圆性是 Hodge 定理成立的分析基础：椭圆 + 自伴 + 非负 + 紧致流形 $\Rightarrow$ 谱离散且有限维核。在物理中，紧致空间上经过适当规范固定的椭圆线性化具有有限维核；这不能直接推广为全部非线性规范场解空间有限维，以及弦理论中 拉普拉斯算子 谱决定物理态能级。Weitzenböck 公式把椭圆性 + 曲率正性结合，给出 Bochner 消没定理——曲率通过低阶项直接控制核空间维数。
\end{proof}

\begin{theorem}[Green 算子与 Hodge 分解]
在紧致、定向、无边界 Riemann 流形上，每个光滑 $k$-形式唯一分解为三个两两正交分量\[\alpha=\Pi_{\mathcal H,k}\alpha+\dd\delta\mathcal G_k\alpha+\delta\dd\mathcal G_k\alpha.\]每个 de Rham 类有唯一调和代表，两个势本身可不唯一。
\end{theorem}

\begin{proof}
在紧致无边界流形上，限制在 $\Omega^k(M)$ 的椭圆自伴算子记为 $\Delta_k$。它具有离散非负谱，零特征空间就是有限维的 $\mathcal H^k(M)$。记到该零空间的正交投影为 $\Pi_{\mathcal H,k}$，并在其正交补上定义 Green 算子
\[
\mathcal G_k:=\left(\Delta_k\big|_{\mathcal H^k(M)^\perp}\right)^{-1},
\qquad
\mathcal G_k\eta=0\quad(\eta\in\mathcal H^k(M)).
\]
于是谱分解给出精确算子恒等式
\[
I=\Pi_{\mathcal H,k}+\Delta_k\mathcal G_k.
\]
对任意 $\alpha\in\Omega^k(M)$，直接展开 $\Delta_k=\dd\delta+\delta\dd$ 得
\[
\begin{aligned}
\alpha
&=\Pi_{\mathcal H,k}\alpha+\Delta_k\mathcal G_k\alpha\\
&=\Pi_{\mathcal H,k}\alpha
+\dd\delta\mathcal G_k\alpha
+\delta\dd\mathcal G_k\alpha.
\end{aligned}
\]
令
\[
\alpha_{\mathcal H}=\Pi_{\mathcal H,k}\alpha,
\qquad
\beta=\delta\mathcal G_k\alpha,
\qquad
\gamma=\dd\mathcal G_k\alpha,
\]
便得到\eqnrefs{eq:math-dg-hodge-decomposition}。

三个分量两两正交。例如
\[
(\dd\beta,\delta\gamma)
=(\dd^2\beta,\gamma)=0,
\]
而对 调和 分量 $\alpha_{\mathcal H}$，由 $\dd\alpha_{\mathcal H}=\delta\alpha_{\mathcal H}=0$，
\[
(\dd\beta,\alpha_{\mathcal H})=(\beta,\delta\alpha_{\mathcal H})=0,
\qquad
(\delta\gamma,\alpha_{\mathcal H})=(\gamma,\dd\alpha_{\mathcal H})=0.
\]
因此三项本身唯一。若只改变 $\beta$ 为 $\beta+\zeta$ 且 $\dd\zeta=0$，精确 分量不变，所以势并不唯一；$\gamma$ 同理。

最后证明调和代表与 de Rham 上同调一一对应。调和形式由\eqnrefs{eq:math-dg-harmonic-closed-coclosed}必为闭形式，故定义上同调类。若 $\eta_1-\eta_2=\dd\zeta$ 且 $\eta_1,\eta_2$ 都 调和，则
\[
\|\eta_1-\eta_2\|^2
=(\eta_1-\eta_2,\dd\zeta)
=(\delta(\eta_1-\eta_2),\zeta)=0,
\]
所以同一上同调类至多一个 调和 代表。另一方面若 $\alpha$ 闭，Hodge 分解给出
\[
\alpha=\dd\beta+\delta\gamma+\alpha_{\mathcal H}.
\]
对两边作用 $\dd$，得到 $\dd\delta\gamma=0$。于是
\[
\|\delta\gamma\|^2
=(\gamma,\dd\delta\gamma)=0,
\]
故 $\delta\gamma=0$，从而 $[\alpha]=[\alpha_{\mathcal H}]$。存在性成立。
\end{proof}

\begin{theorem}[实 de Rham 上同调的 Poincaré 对偶]
设 $M^n$ 紧致、定向、无边界。配对\[H^k_{\mathrm{dR}}(M;\RR)\times H^{n-k}_{\mathrm{dR}}(M;\RR)\to\RR,\qquad([\alpha],[\beta])\mapsto\int_M\alpha\wedge\beta\]良定义且非退化，故 $b_k=b_{n-k}$。
\end{theorem}

\begin{proof}
先验证\eqnrefs{eq:math-dg-poincare-pairing}只依赖上同调类。若 $\alpha\mapsto\alpha+\dd\eta$ 且 $\beta$ 闭，则由 $\dd\beta=0$ 及外微分 Leibniz 律
\[
\dd(\eta\wedge\beta)
=\dd\eta\wedge\beta+(-1)^{\deg\eta}\eta\wedge\dd\beta
=\dd\eta\wedge\beta,
\]
故由 Stokes 定理及 $\partial M=\varnothing$，
\[
\int_M\dd\eta\wedge\beta
=\int_M\dd(\eta\wedge\beta)
=\int_{\partial M}\eta\wedge\beta
=0.
\]
对 $\beta$ 的 精确 改变 $\beta\mapsto\beta+\dd\zeta$ 同理，因 $\alpha$ 闭且
\[
\dd(\alpha\wedge\zeta)
=\dd\alpha\wedge\zeta+(-1)^{\deg\alpha}\alpha\wedge\dd\zeta
=(-1)^{\deg\alpha}\alpha\wedge\dd\zeta,
\]
故
\[
\int_M\alpha\wedge\dd\zeta
=(-1)^k\int_M\dd(\alpha\wedge\zeta)=0.
\]
因此配对在上同调上良定。

现在取非零上同调类 $[\alpha]\neq0$，并令 $h$ 为其唯一调和代表元，即 $\Delta h=0$。由\eqnrefs{eq:math-dg-hodge-star-harmonic-isomorphism}，Hodge 星保持 调和 性，故 $*h$ 也是 调和，因此定义一个 $(n-k)$ 次上同调类
\[
[*h]\in H^{n-k}_{\mathrm{dR}}(M).
\]
于是由 Hodge 星的定义\eqnrefs{eq:math-dg-hodge-star-characterization}，
\begin{align*}
\int_M h\wedge *h
&=\int_M \langle h,h\rangle_g\,\operatorname{vol}_g\\
&=\int_M |h|_g^2\,\operatorname{vol}_g\\
&=\|h\|_{L^2}^2.
\end{align*}
由 $[\alpha]\neq0$ 故 $h\neq0$，从而
\[
\|h\|_{L^2}^2>0.
\]
所以 $[\alpha]$ 与 $[*h]$ 配对非零，不可能与所有 $H^{n-k}_{\mathrm{dR}}(M)$ 元素配对为零。反向非退化性由 $**h=(-1)^{k(n-k)}h$ 完全相同地得到：对非零 $[\beta]\in H^{n-k}_{\mathrm{dR}}(M)$，取其 调和 代表元 $k$，则
\[
\int_M k\wedge *k=\|k\|_{L^2}^2>0,
\]
故 $[\beta]$ 与 $[*k]\in H^k_{\mathrm{dR}}(M)$ 配对非零。因此配对非退化，而有限维向量空间之间的非退化配对立刻给出
\[
\dim H^k_{\mathrm{dR}}(M)
=\dim H^{n-k}_{\mathrm{dR}}(M).
\]

紧曲面 $\Sigma_g$（亏格 $g$）满足 $b_0=b_2=1$、$b_1=2g$，因而 $b_k=b_{2-k}$ 逐次成立；中间次数 $k=1$ 是自对偶的。紧 3-流形如 $T^3$ 则有 $b_0=b_3=1$、$b_1=b_2=3$。Poincaré 对偶不仅是维数等式，还给出上积配对 $H^k\times H^{n-k}\to H^n\cong\RR$ 的非退化性——这是相交理论的基础：$k$-形式与 $(n-k)$-形式的上积对应对偶子流形的相交数。
\end{proof}

\begin{proposition}[一形式的 Weitzenböck 公式]
对 Riemann 流形上的光滑一形式，\[(\Delta_H\alpha)_j=(\nabla^*\nabla\alpha)_j+\operatorname{Ric}_j{}^k\alpha_k.\]在紧致无边界且 $\operatorname{Ric}>0$ 时，每个调和一形式为零。
\end{proposition}

\begin{proof}
对 $1$-形式 $\alpha$，分别计算 $\delta\dd\alpha$ 与 $\dd\delta\alpha$ 的分量，相加后把"协变导数交换子"识别为 Ricci 曲率，得到 $\Delta_H=\nabla^*\nabla+\operatorname{Ric}$。

对 $1$-形式 $\alpha=\alpha_j\dd x^j$，采用
\begin{equation*}
(\dd\alpha)_{ij}
=\nabla_i\alpha_j-\nabla_j\alpha_i,
\qquad
\delta\alpha=-\nabla^i\alpha_i.
\end{equation*}
第一个是外微分的协变分量形式（无挠使 Christoffel 项相消），第二个是余微分的分量约定（散度的负号来自本书非负 拉普拉斯算子 约定）。

把 $\dd\alpha$ 视为 $2$-形式再作用 $\delta$：
\begin{equation*}
(\delta\dd\alpha)_j
=-\nabla^i(\dd\alpha)_{ij}
=-\nabla^i(\nabla_i\alpha_j-\nabla_j\alpha_i)
=-\nabla^i\nabla_i\alpha_j+\nabla^i\nabla_j\alpha_i.
\end{equation*}
第一项是非负约定的联络拉普拉斯 $\nabla^*\nabla$ 的主项，第二项保留待与 $\dd\delta\alpha$ 中的对应项合并。


\begin{equation*}
(\dd\delta\alpha)_j
=\nabla_j(\delta\alpha)
=-\nabla_j\nabla^i\alpha_i.
\end{equation*}
此处用了 $1$-形式 $\delta\alpha$ 的外微分分量公式 $(\dd f)_j=\nabla_j f$。

相加得到
\begin{equation*}
(\Delta_H\alpha)_j
=(\delta\dd\alpha+\dd\delta\alpha)_j
=-\nabla^i\nabla_i\alpha_j
+\nabla^i\nabla_j\alpha_i
-\nabla_j\nabla^i\alpha_i.
\end{equation*}
最后两项正是协变导数交换子。对协变 $1$-张量，
\begin{equation*}
(\nabla^i\nabla_j-\nabla_j\nabla^i)\alpha_i
=\operatorname{Ric}_j{}^k\alpha_k
\end{equation*}
（符号与本书前面的 Riemann 曲率 约定 相容）。这一识别来自 Riemann 曲率的分量定义 $[\nabla_a,\nabla_b]X^c=R^c{}_{dab}X^d$，对协变指标取 迹 给 Ricci：$R^i{}_{j i k}\alpha^k=\operatorname{Ric}_{jk}\alpha^k$。

另一方面，正规坐标中
\begin{equation*}
(\nabla^*\nabla\alpha)_j
=-\nabla^i\nabla_i\alpha_j
\end{equation*}
在基点成立，张量性使其全局成立，因此得到\eqnrefs{eq:math-dg-weitzenbock-oneform}。

典型例子与应用。
\begin{itemize}
\item $\RR^n$：$\operatorname{Ric}=0$，$\Delta_H=\nabla^*\nabla=-\sum_i\partial_i^2$，每个分量独立满足标量 拉普拉斯算子。
\item $S^n$（$n\ge2$，Ric $=(n-1)g$）：$\Delta_H=\nabla^*\nabla+(n-1)$。调和 $1$-形式满足 $\nabla^*\nabla\alpha=-(n-1)\alpha$，但 $\nabla^*\nabla$ 非负，故 $\alpha=0$，即 $H^1(S^n)=0$（$n\ge2$）。
\item 紧致流形 $\operatorname{Ric}>0$：Bochner 消没定理给 $H^1(M)=0$，结合 Bonnet--Myers 定理（$\operatorname{Ric}\ge(n-1)k>0$ $\Rightarrow$ 直径 $\le\pi/\sqrt{k}$）说明正 Ricci 曲率强烈限制流形拓扑。
\end{itemize}

对紧致、无边界的 Spin 流形，未扭曲 Dirac 算子满足
$D^2=\nabla^*\nabla+\mathrm{Scal}/4$。
若 $\mathrm{Scal}>0$，能量恒等式推出 $\ker D=0$。
在偶数维，进一步结合分次算子 $D^+$ 的 Atiyah--Singer 指标公式，
得到 $\operatorname{ind}(D^+)=\hat A(M)=0$。
若加入非平坦扭曲丛，平方中还有丛曲率项，不能只凭正标量曲率断言核为零。
热核的分次迹 $\operatorname{Str}(\ee^{-tD^2})$ 将局部曲率项与全局指标联系起来，
这一联系将在旋量章节的 Lichnerowicz 公式中继续展开。
\end{proof}

\begin{theorem}[Lichnerowicz 谱隙下界]
设 $(M^n,g)$ 紧致、连通、无边界，$n\ge2$，$\operatorname{Ric}\ge(n-1)kg$、$k>0$。非负函数拉普拉斯 $\Delta=-\Delta_{\mathrm{LB}}$ 的第一正本征值满足 $\lambda_1\ge nk$。
\end{theorem}

\begin{proof}
设 $(M^n,g)$ 紧致无边界，$\operatorname{Ric}\ge(n-1)kg$（$k>0$），$\Delta f=\lambda_1 f$（第一非零本征函数）。

对本征函数积分\eqnrefs{eq:math-dg-bochner-function}。紧致无边界时 Stokes 定理给出 $\int_M\Delta|\nabla f|^2\operatorname{vol}_g=0$（边界项消失），于是
\begin{equation*}
 0
 =\lambda_1\int_M|\nabla f|^2\operatorname{vol}_g
 -\int_M|\nabla^2f|^2\operatorname{vol}_g
 -\int_M\operatorname{Ric}(\nabla f,\nabla f)\operatorname{vol}_g.
\end{equation*}

由 Green 公式，
\begin{equation*}
 \int_M|\nabla f|^2\operatorname{vol}_g
 =\int_M f\,\Delta f\,\operatorname{vol}_g
 =\lambda_1\int_M f^2\operatorname{vol}_g.
\end{equation*}

由 Cauchy--Schwarz 不等式对 Hessian 的迹，
\begin{equation*}
 |\nabla^2 f|^2\ge\frac{1}{n}(\tr\nabla^2f)^2
 =\frac{1}{n}(\Delta f)^2
 =\frac{\lambda_1^2}{n}f^2.
\end{equation*}
由 Ricci 下界假设，
\begin{equation*}
 \operatorname{Ric}(\nabla f,\nabla f)\ge(n-1)k\,|\nabla f|^2.
\end{equation*}

这里只能在积分以后使用
$\int|\nabla f|^2=\lambda_1\int f^2$，点态恒等式 $|\nabla f|^2=\lambda_1f^2$ 不成立。
将两个下界代入上面的 Bochner 恒等式（移项后右端 $\ge0$）：
\begin{align*}
 0
 &\le \lambda_1^2\int_Mf^2\operatorname{vol}_g
  -\frac{\lambda_1^2}{n}\int_Mf^2\operatorname{vol}_g
  -(n-1)k\lambda_1\int_Mf^2\operatorname{vol}_g\\
 &=\lambda_1\left(\frac{(n-1)\lambda_1}{n}-(n-1)k\right)\int_Mf^2\operatorname{vol}_g\\
 &=(n-1)\lambda_1\left(\frac{\lambda_1}{n}-k\right)\int_Mf^2\operatorname{vol}_g.
\end{align*}

因为 $\lambda_1>0$（非常数本征函数）、$n>1$、$f\not\equiv0$，积分因子为正，故 $\lambda_1/n-k\ge0$，即 $\lambda_1\ge nk$。这就是\eqnrefs{eq:math-dg-lichnerowicz-gap}。

Lichnerowicz 下界与 Bonnet--Myers 直径上界 $\operatorname{diam}(M)\le\pi/\sqrt{k}$ 假设完全相同（$\operatorname{Ric}\ge(n-1)kg$），但结论从"直径有限"变成"拉普拉斯算子 的第一个正本征值远离零"。同一个 Ricci 下界同时约束几何尺度与低频谱。
\end{proof}

\begin{proposition}[Hodge 边界条件的 Green 能量恒等式]
设 $M$ 紧致、定向且有光滑边界。绝对条件 $\iota_\nu\omega=\iota_\nu\dd\omega=0$ 或相对条件 $j^*\omega=j^*\delta\omega=0$ 均推出\[(\Delta_H\omega,\omega)=\|\dd\omega\|^2+\|\delta\omega\|^2.\]调和零模闭且余闭，Hodge 星给出 $\mathcal H_A^k\cong\mathcal H_R^{n-k}$。
\end{proposition}

\begin{proof}
设 $M$ 为紧致、定向的 Riemann 流形，边界光滑，$\nu$ 为外法单位向量。
先对实微分形式计算；复形式在第二槽取共轭，结论相同。
对 $\alpha\in\Omega^{k-1}(M)$、$\beta\in\Omega^k(M)$，外微分乘积规则给出
\[
\dd(\alpha\wedge *\beta)
=\dd\alpha\wedge *\beta+(-1)^{k-1}\alpha\wedge\dd(*\beta).
\]
按正文的余微分约定，第二项积分为 $-(\alpha,\delta\beta)$。
Stokes 定理因而给出
\[
(\dd\alpha,\beta)-(\alpha,\delta\beta)
=\int_{\partial M}j^*(\alpha\wedge *\beta)
=\int_{\partial M}\langle\alpha,\iota_\nu\beta\rangle\,\dd S.
\]
最后一个等号可在正交余标架 $(\nu^\flat,\theta^1,\dots)$ 中验证：
只有 $\beta$ 的法向分量与 $\alpha$ 的切向分量参与边界积分。

对 $\Delta\omega=\delta\dd\omega+\dd\delta\omega$ 分别使用上式，得
\begin{align*}
(\delta\dd\omega,\omega)
&=\|\dd\omega\|^2-
\int_{\partial M}\langle\omega,\iota_\nu\dd\omega\rangle\,\dd S,\\
(\dd\delta\omega,\omega)
&=\|\delta\omega\|^2+
\int_{\partial M}\langle\delta\omega,\iota_\nu\omega\rangle\,\dd S.
\end{align*}
绝对条件 $\iota_\nu\omega=0$、$\iota_\nu\dd\omega=0$
逐项消去边界积分。相对条件 $j^*\omega=0$、$j^*\delta\omega=0$
则消去其中的切向因子。故两种域都满足
\[
(\Delta\omega,\omega)=\|\dd\omega\|^2+\|\delta\omega\|^2.
\]
若 $\Delta\omega=0$，两项非负，必有 $\dd\omega=\delta\omega=0$。
能量恒等式证明非负性和零模的闭、余闭性质；完整自伴性还使用
Hodge--Morrey--Friedrichs 的椭圆边界理论，不能只由一次分部积分宣称已经证明。

Hodge 星的点态恒等式
\[
\iota_\nu(*\omega)=*(\omega\wedge\nu^\flat),\qquad
**\omega=(-1)^{k(n-k)}\omega
\]
说明它交换切向和法向分量。再利用 $\delta*=\pm*\dd$，
绝对条件的两条式子变为 $j^*(*\omega)=0$、$j^*\delta(*\omega)=0$。
因此 $*:\mathcal H_A^k\to\mathcal H_R^{n-k}$ 是同构。

区间上的 $0$-形式使两种条件完全可见：绝对条件为
$f'(0)=f'(1)=0$，相对条件为 $f(0)=f(1)=0$。
$-f''=0$ 的 Neumann 解是常数，而 Dirichlet 解只有零。
相应地，$H^0([0,1])=\RR$、$H^0([0,1],\partial[0,1])=0$。

在乘积平坦度量的圆柱 $S^1\times[0,1]$ 上，$\dd\theta$ 是绝对调和
$1$-形式，$\dd t$ 是相对调和 $1$-形式。两种空间维数都为一，
但选出的代表不同；维数不能被写成必然不同。
圆盘以及半球同胚于闭球 $D^n$，其上同调为
\[
H^k(D^n)=
\begin{cases}\RR,&k=0,\\0,&k>0,\end{cases}
\qquad
H^k(D^n,\partial D^n)=
\begin{cases}\RR,&k=n,\\0,&k\ne n.\end{cases}
\]
特别地 $H^1(D^2,\partial D^2)=0$，非平凡相对类在次数 $2$，
由体积形式代表。半球的普通上同调不能等同于整个球面 $S^n$ 的上同调。
这些结果也可由可缩性、对偶定理与配对长正合列分别核对。

无边界时两种域重合，恢复通常的 Hodge 分解。
非定向的 Möbius 带则需要取向线丛：Hodge 星把普通形式映为取向扭曲的形式，
也可以在可定向二重覆盖上施加正确的奇偶条件，不能忽略该变换的类型。
\end{proof}
